\documentclass[handout]{beamer}
\mode<presentation>
\usepackage{xcolor}
\usepackage[toc]{multitoc}
\usepackage{color}
\usepackage{graphicx}
\graphicspath{{../buku_ajar/buku_ajar_logika_matematika/figures/}}
\usepackage{tikz}
\usetikzlibrary{shapes,arrows,positioning,calc}
\usepackage{listings}
\usepackage{lmodern}
\usepackage{amsmath}
\usetheme[secheader]{Boadilla}
\usecolortheme{whale}
\setbeamertemplate{caption}[numbered]
\setbeamertemplate{bibliography item}[text]
\setbeamertemplate{navigation symbols}{}

\newcommand{\logic}[1]{\textbf{\textit{#1}}}

\theoremstyle{definition}
\newtheorem{definisi}[theorem]{Definisi}
\theoremstyle{example}
\newtheorem{contoh}[theorem]{Contoh}

\renewcommand*{\multicolumntoc}{2}

\setbeamertemplate{footline}
{
  \leavevmode%
  \hbox{%
  \begin{beamercolorbox}[wd=.333333\paperwidth,ht=2.25ex,dp=1ex,center]{author in head/foot}%
    \usebeamerfont{author in head/foot}\insertshortauthor%
  \end{beamercolorbox}%
  \begin{beamercolorbox}[wd=.433333\paperwidth,ht=2.25ex,dp=1ex,center]{title in head/foot}%
    \usebeamerfont{title in head/foot}\insertshorttitle
  \end{beamercolorbox}%
  \begin{beamercolorbox}[wd=.233333\paperwidth,ht=2.25ex,dp=1ex,right]{date in head/foot}%
    \usebeamerfont{date in head/foot}\insertshortdate{}\hspace*{2em} \insertframenumber{} / \inserttotalframenumber\hspace*{2ex} 
  \end{beamercolorbox}}%
  \vskip0pt%
}

\subtitle[Logika Matematika]{Logika Matematika}
\title[Induksi Matematika]{Prinsip Induksi Matematika}
\author{Cecep Suwanda, S.Si., M.Kom.}
\date{Maret 2025}
\institute{Universitas Bale Bandung}

\begin{document}

\maketitle

\section*{Daftar Isi}
\begin{frame}
\frametitle{Daftar Isi}
\tableofcontents
\end{frame}

% ============================================================
% SECTION 13-1: Prinsip dan domino
% ============================================================
\section{Prinsip Induksi}

\begin{frame}
\frametitle{Kapan Induksi Digunakan?}
\begin{itemize}
  \item Pernyataan untuk bilangan bulat berurutan ($n \geq n_0$)
  \item Deret, formula penjumlahan, pertidaksamaan, keterbagian
  \item Domain tak berhingga (bukan verifikasi exhaustif)
\end{itemize}
\end{frame}

\begin{frame}
\frametitle{Analogi Efek Domino}
\begin{enumerate}
  \item \textbf{Basis}: Batu pertama jatuh ($P(n_0)$ benar)
  \item \textbf{Induksi}: Jika batu ke-$k$ jatuh, batu ke-$(k+1)$ jatuh
\end{enumerate}
Kedua syarat wajib; tanpa basis, rantai tidak dimulai.
\end{frame}

\begin{frame}
\frametitle{Prinsip Induksi Matematika}
\begin{definisi}[Prinsip Induksi]
Jika (1) $P(n_0)$ benar \textbf{(basis)}, dan\\
(2) $\forall k \geq n_0$: $P(k) \Rightarrow P(k+1)$ \textbf{(langkah induksi)},\\
maka $P(n)$ benar untuk semua $n \geq n_0$.
\end{definisi}
\[
\bigl[ P(n_0) \wedge \forall k \geq n_0\, (P(k) \to P(k+1)) \bigr] \to \forall n \geq n_0\, P(n)
\]
$n_0$ tidak harus 1 (bisa 0, 2, \ldots).
\end{frame}

% ============================================================
% SECTION 13-2: Induksi lemah / template
% ============================================================
\section{Induksi Lemah}

\begin{frame}
\frametitle{Template Bukti Induksi}
\begin{enumerate}
  \item Nyatakan $P(n)$ yang dibuktikan
  \item \textbf{Basis}: verifikasi $P(n_0)$
  \item \textbf{Hipotesis induksi}: asumsikan $P(k)$ untuk $k \geq n_0$
  \item \textbf{Langkah induksi}: tunjukkan $P(k+1)$ memakai $P(k)$
  \item Kesimpulan: menurut induksi, $P(n)$ untuk semua $n \geq n_0$
\end{enumerate}
\end{frame}

\begin{frame}
\frametitle{Contoh: Jumlah $n$ Bilangan Ganjil Pertama}
$1 + 3 + 5 + \cdots + (2n-1) = n^2$\\
Basis $n=1$: $1 = 1^2$.\\
Induksi: $k^2 + (2k+1) = (k+1)^2$.
\end{frame}

\begin{frame}
\frametitle{Pentingnya Langkah Basis}
Tanpa basis, rantai implikasi tidak pernah ``dipicu''.\\
Induksi tanpa basis = domino tanpa dorongan pertama (tidak valid).
\end{frame}

% ============================================================
% SECTION 13-3: Deret
% ============================================================
\section{Induksi pada Deret}

\begin{frame}
\frametitle{Jumlah Bilangan 1 sampai $n$}
$\displaystyle\sum_{i=1}^{n} i = \frac{n(n+1)}{2}$\\
Pola induksi: pisahkan suku $(k+1)$, pakai hipotesis pada $1+\cdots+k$, sederhanakan ke bentuk formula untuk $n=k+1$.
\end{frame}

\begin{frame}
\frametitle{Jumlah Kuadrat dan Deret Geometri}
\begin{itemize}
  \item $\displaystyle\sum_{i=1}^{n} i^2 = \frac{n(n+1)(2n+1)}{6}$
  \item $\displaystyle\sum_{i=0}^{n} r^i = \frac{r^{n+1}-1}{r-1}$ ($r \neq 1$, basis $n=0$)
\end{itemize}
\end{frame}

% ============================================================
% SECTION 13-4: Pertidaksamaan
% ============================================================
\section{Induksi pada Pertidaksamaan}

\begin{frame}
\frametitle{Contoh: $2^n > n$ untuk $n \geq 1$}
Basis: $2^1 > 1$.\\
Induksi: $2^{k+1} = 2 \cdot 2^k > 2k \geq k+1$ (untuk $k \geq 1$).
\end{frame}

\begin{frame}
\frametitle{Basis Sesuai Domain}
Contoh: $n! > 2^n$ untuk $n \geq 4$ --- basis di $n=4$, bukan $n=1$.\\
Bernoulli: $(1+x)^n \geq 1+nx$ untuk $x \geq -1$, $n \geq 1$.
\end{frame}

% ============================================================
% SECTION 13-5: Keterbagian
% ============================================================
\section{Induksi pada Keterbagian}

\begin{frame}
\frametitle{Contoh: $3 \mid (n^3 - n)$}
Basis: $1^3-1=0$.\\
Induksi: $(k+1)^3-(k+1) = (k^3-k) + 3(k^2+k)$ = kelipatan 3.
\end{frame}

\begin{frame}
\frametitle{Contoh Lain}
\begin{itemize}
  \item $6 \mid n(n+1)(n+2)$ untuk $n \geq 1$
  \item $3 \mid (4^n - 1)$ untuk $n \geq 0$
\end{itemize}
Teknik: pemfaktoran + substitusi hipotesis induksi.
\end{frame}

% ============================================================
% SECTION 13-6: Induksi kuat
% ============================================================
\section{Induksi Kuat}

\begin{frame}
\frametitle{Prinsip Induksi Kuat}
\begin{definisi}[Induksi Kuat]
(1) $P(n_0)$ benar;\\
(2) Untuk $k \geq n_0$: jika $P(n_0),\ldots,P(k)$ benar maka $P(k+1)$ benar.\\
$\Rightarrow$ $P(n)$ benar untuk semua $n \geq n_0$.
\end{definisi}
Hipotesis: \textbf{semua} $P(i)$ dari $n_0$ sampai $k$, bukan hanya $P(k)$.
\end{frame}

\begin{frame}
\frametitle{Lemah vs Kuat}
\begin{center}
\small
\begin{tabular}{|l|l|l|}
\hline
 & \textbf{Lemah} & \textbf{Kuat} \\ \hline
Asumsi & $P(k)$ & $P(n_0),\ldots,P(k)$ \\ \hline
Kapan & $P(k+1)$ dari $P(k)$ saja & Bergantung beberapa nilai sebelumnya \\ \hline
Contoh & Deret, pertidaksamaan & Faktor prima, Fibonacci \\ \hline
\end{tabular}
\end{center}
\end{frame}

\begin{frame}
\frametitle{Faktorisasi Prima}
Setiap $n \geq 2$ = hasil kali prima.\\
Jika $k+1$ komposit: $k+1 = ab$ dengan $a,b \leq k$ --- perlu $P(a), P(b)$ (induksi kuat).
\end{frame}

\begin{frame}
\frametitle{Fibonacci: $F_n < 2^n$}
Dua basis: $F_1 < 2^1$, $F_2 < 2^2$.\\
$F_{k+1} = F_k + F_{k-1} < 2^k + 2^{k-1} < 2^{k+1}$.
\end{frame}

% ============================================================
% SECTION 13-7: Kesalahan umum
% ============================================================
\section{Kesalahan Umum}

\begin{frame}
\frametitle{Kesalahan yang Sering Terjadi}
\begin{enumerate}
  \item Melupakan langkah basis
  \item Tidak memakai hipotesis induksi pada langkah induksi
  \item Arah salah / sirkuler (mulai dari yang ingin dibuktikan)
  \item Basis salah (mis.\ $n \geq 4$ tetapi basis di $n=1$)
\end{enumerate}
\end{frame}

\begin{frame}
\frametitle{Panduan Praktis}
Identifikasi $P(n)$ dan $n_0$ $\to$ verifikasi basis $\to$ tulis hipotesis $\to$ tujuan $P(k+1)$ $\to$ kerjakan dari kiri, substitusi $P(k)$ $\to$ simpulkan.
\end{frame}

\begin{frame}
\frametitle{Hubungan ke Bab Berikutnya}
Bab 14: \textbf{loop invariant} --- analogi basis (kondisi awal) dan langkah induksi (pemeliharaan tiap iterasi).
\end{frame}

% ============================================================
% RANGKUMAN
% ============================================================
\section{Rangkuman}

\begin{frame}
\frametitle{Rangkuman}
\begin{itemize}
  \item Basis \textbf{dan} langkah induksi keduanya wajib
  \item Induksi lemah: asumsikan $P(k)$, buktikan $P(k+1)$
  \item Induksi kuat: asumsikan $P(n_0),\ldots,P(k)$
  \item Aplikasi: deret, pertidaksamaan, keterbagian, rekursi
  \item Hindari basis terlewat, hipotesis tidak dipakai, basis salah
\end{itemize}
\end{frame}

\begin{frame}
\frametitle{Kata Kunci}
\begin{center}
\small
induksi matematika $\bullet$ langkah basis $\bullet$ langkah induksi $\bullet$ hipotesis induksi $\bullet$ induksi lemah $\bullet$ induksi kuat $\bullet$ efek domino $\bullet$ deret $\bullet$ pertidaksamaan $\bullet$ keterbagian $\bullet$ Fibonacci
\end{center}
\end{frame}

\end{document}
